\newcommand{\showviz}[2]{%
  \resizebox{#1\linewidth}{!}{%
	\setlength{\fboxsep}{0pt}
		\includegraphics[height=#1\linewidth]{images/#2.png}
		\hfill
		\includegraphics[height=#1\linewidth]{images/#2_gt.png}
		\hfill
		\includegraphics[height=#1\linewidth]{images/#2_v2.png}
		\hfill
		\includegraphics[height=#1\linewidth]{images/#2_v1.png}
        \hfill
		\includegraphics[height=#1\linewidth]{images/#2_siglip2.png}
}}


\newcommand{\vizwidth}{1.0}
\begin{figure}[tb]
	\centering
	\showviz{\vizwidth}{dog_rug} \\[0.2mm]
    \showviz{\vizwidth}{birds} \\[0.2mm]
    \showviz{\vizwidth}{bus} \\[0.2mm]

    \resizebox{\vizwidth\linewidth}{!}{%
    \scriptsize
        \begin{minipage}[t]{0.19\linewidth}
          \centering
          Image
        \end{minipage}%
        \hfill
        \begin{minipage}[t]{0.19\linewidth}
          \centering
          Ground truth
        \end{minipage}%
        \hfill
        \begin{minipage}[t]{0.19\linewidth}
          \centering
          \tipsvtwo (ours)
        \end{minipage}%
        \hfill
        \begin{minipage}[t]{0.19\linewidth}
          \centering
          TIPS
        \end{minipage}%
        \hfill
        \begin{minipage}[t]{0.19\linewidth}
          \centering
          SigLIP2
        \end{minipage}%
    }%
    \vspace{-1mm}
	\caption{\small\textbf{Zero-shot segmentation visualization.} Comparing results for \tipsvtwo, TIPS and SigLIP 2, where classes are predicted directly by finding the closest text embedding to each image patch token, without any post-processing. \tipsvtwo enhances patch-text alignment significantly compared to other models, showcasing strong off-the-shelf capabilities.}
	\label{fig:zss_viz}
	\vspace{-5mm}
\end{figure}

\section{Conclusion}

In this work, we introduce \tipsvtwo, a new vision-language encoder suitable to a variety of multimodal applications.
By investigating the weak dense image-text alignment of previous models, we find that patch-level distillation substantially enhances this capability.
This observation inspires us to introduce iBOT++, an upgrade to the well-known iBOT masked image modeling technique used in pretraining, resulting in significant enhancements to dense image-text alignment.
Additionally, we make vision-language pretraining more efficient with simplifications to the EMA setup, and leverage image captions at different granularities to obtain more robust representations.
Through comprehensive experiments, our final \tipsvtwo model showcases strong results on $9$ tasks spanning $20$ datasets, generally matching or exceeding results from previous vision-language encoders.

\noindent\textbf{Acknowledgments.}~We would like to thank Connor Schenck for thoughtful discussions and suggestions.
